\section{Непрерывность функции в точке, на отрезке, на интервале. Область
непрерывности функций: степенной (для натурального показателя),
тригонометрических.}

\subsection{Непрерывность функции в точке}

\begin{definition}
Функция $f$ называется непрерывной в точке $x_0$, если:
\begin{enumerate}
\item $f(x_0)$ определена;
\item существует предел $\Lim{x}{x_0} f(x)$;
\item $\Lim{x}{x_0} f(x)=f(x_0)$.
\end{enumerate}
\end{definition}

Эквивалентная запись:
\[
\Lim{x}{x_0}\bigl(f(x)-f(x_0)\bigr)=0.
\]

В $\varepsilon$–$\delta$ форме:
\[
\forall \varepsilon>0\ \exists \delta>0:
|x-x_0|<\delta \Rightarrow |f(x)-f(x_0)|<\varepsilon.
\]

\subsection{Непрерывность на множестве, отрезке и интервале}

\begin{definition}
Функция называется непрерывной на множестве $E$, если она непрерывна в каждой точке этого множества.
\end{definition}

Для отрезка $[a,b]$:
\begin{itemize}
\item $[a;b] \subset D_f$
\item в точках $(a,b)$ — обычная непрерывность (непрерывна на интервале --- в каждой его точке);
\item в точке $a$ — непрерывность справа;
\item в точке $b$ — непрерывность слева.
\end{itemize}

\subsection{Свойства непрерывных функций}

\begin{theorem}
Пусть \(f\) и \(g\) непрерывны в точке \(x_0\). Тогда в точке \(x_0\) непрерывны функции
\[
f+g,\qquad f-g,\qquad f\cdot g,
\]
а также частное
\[
\frac{f}{g},
\]
если \(g(x_0)\neq0\).
\end{theorem}

\begin{Proof}
Так как \(f\) и \(g\) непрерывны в \(x_0\), то
\[
\lim_{x\to x_0} f(x)=f(x_0),\qquad \lim_{x\to x_0} g(x)=g(x_0).
\]

Пользуясь свойствами пределов получаем
\[
\lim_{x\to x_0}\bigl(f(x)+g(x)\bigr)
=\lim_{x\to x_0}f(x)+\lim_{x\to x_0}g(x)
=f(x_0)+g(x_0),
\]
откуда следует непрерывность \(f+g\) в \(x_0\).

Аналогично,
\[
\lim_{x\to x_0}\bigl(f(x)-g(x)\bigr)
=f(x_0)-g(x_0),
\]
и
\[
\lim_{x\to x_0}\bigl(f(x)\,g(x)\bigr)
=\bigl(\lim_{x\to x_0}f(x)\bigr)\cdot\bigl(\lim_{x\to x_0}g(x)\bigr)
=f(x_0)\,g(x_0),
\]
что даёт непрерывность разности и произведения.

Для частного: если \(g(x_0)\neq0\), то из непрерывности \(g\) следует, что \(\lim_{x\to x_0}g(x)=g(x_0)\neq0\). Следовательно, в некоторой окрестности точки \(x_0\) функция \(g(x)\) не обращается в ноль, и можно применять правило для предела частного:
\[
\lim_{x\to x_0}\frac{f(x)}{g(x)}
=\frac{\lim\limits_{x\to x_0}f(x)}{\lim\limits_{x\to x_0}g(x)}
=\frac{f(x_0)}{g(x_0)}.
\]
Отсюда \(\dfrac{f}{g}\) непрерывна в \(x_0\).
\end{Proof}

\subsection{Связь с дифференцируемостью}

Если функция дифференцируема в точке, то она непрерывна в этой точке.

\begin{Proof}
По определению дифференцируемости:
\[
f'(x_0) = \Lim{x}{x_0} \frac{f(x)-f(x_0)}{x-x_0}
\]
существует. Умножая обе части на $(x-x_0)$ и беря предел, получаем
\[
\Lim{x}{x_0} \bigl(f(x)-f(x_0)\bigr) = 0,
\]
то есть $f$ непрерывна в $x_0$.
\end{Proof}

\begin{remark}
Обратное неверно: непрерывность не гарантирует дифференцируемость. 

\textbf{Контрпример:} $f(x)=|x|$ в точке $x_0=0$. 
Функция непрерывна, но $f'(0)$ не существует:
\[
\lim_{x\to 0^+} \frac{|x|-0}{x} = 1, \quad
\lim_{x\to 0^-} \frac{|x|-0}{x} = -1.
\]
\end{remark}

\subsection{Области непрерывности элементарных функций}

\subsubsection{Степенные функции}

\begin{theorem}
Степенная функция
\[
f(x)=x^n, \qquad n\in\mathbb{N},
\]
непрерывна на всей числовой прямой $\mathbb{R}$.
\end{theorem}

\begin{Proof}
При $n=1$ функция $f(x)=x$ непрерывна.

Пусть $n>1$. Тогда
\[
x^n=x\cdot x^{n-1}.
\]

Функция $x$ непрерывна, а по индукционному предположению
$x^{\,n-1}$ также непрерывна. Так как произведение непрерывных
функций непрерывно, то $x^n$ непрерывна на $\mathbb{R}$.
\end{Proof}

\begin{remark}
Если показатель степени $\alpha\in\mathbb{R}$, то функция
\[
f(x)=x^\alpha = e^{\alpha \ln x}
\]
обычно рассматривается на интервале $(0,+\infty)$,
где она непрерывна.
\end{remark}

\subsubsection{Показательная функция}

\begin{theorem}
Функция
\[
f(x)=a^x, \qquad a>0,\ a\neq1,
\]
непрерывна на всей числовой прямой $\mathbb{R}$.
\end{theorem}

\begin{Proof}
Показательная функция определяется через предел степенной
последовательности или через экспоненту $e^{x\ln a}$.
Так как функции $\ln x$ и $e^x$ непрерывны на своих областях
определения, их композиция
\[
a^x=e^{x\ln a}
\]
непрерывна на $\mathbb{R}$.
\end{Proof}

\subsubsection{Тригонометрические функции}

\begin{theorem}
Функции
\[
\sin x, \qquad \cos x
\]
непрерывны на всей $\mathbb{R}$.
\end{theorem}

\begin{Proof}
Из первого замечательного предела следует
\[
\lim_{x\to0}\sin x=0, \qquad
\lim_{x\to0}\cos x=1.
\]

Следовательно, $\sin x$ и $\cos x$ непрерывны в нуле.
Используя формулы сложения
\[
\sin(x+h)=\sin x\cos h+\cos x\sin h,
\]
\[
\cos(x+h)=\cos x\cos h-\sin x\sin h,
\]
получаем непрерывность этих функций в любой точке $x\in\mathbb{R}$.
\end{Proof}

\begin{theorem}
Функции
\[
\tan x=\frac{\sin x}{\cos x}, \qquad
\cot x=\frac{\cos x}{\sin x}
\]
непрерывны во всех точках своей области определения.
\end{theorem}

\begin{Proof}
Функции $\sin x$ и $\cos x$ непрерывны на $\mathbb{R}$.
Частное непрерывных функций непрерывно в точках,
где знаменатель не равен нулю.

Следовательно
\[
\tan x \text{ непрерывна при } x\neq \frac{\pi}{2}+k\pi,
\]

\[
\cot x \text{ непрерывна при } x\neq k\pi,
\quad k\in\mathbb{Z}.
\]
\end{Proof}

\subsubsection{Обратные тригонометрические функции}

\begin{theorem}
Функции
\[
\arcsin x, \qquad \arccos x
\]
непрерывны на интервале
\[
[-1,1].
\]

Функции
\[
\arctan x, \qquad \operatorname{arccot} x
\]
непрерывны на всей $\mathbb{R}$.
\end{theorem}
